\scanpage{143}
ss-gr. \uline{discret} de rang $r_1+r_2\unsure{-}1$, de sorte que le quotient
$V(\mathbb R)/F$ est compact. Il en est à fortiori de même de
$V(\mathbb R)/\overline F$, donc de $V'_{\mathbb R}$ qu'il majore.

\deleted{$(\mathrm{iii})\Longrightarrow(\mathrm{i})$, car si $V'_{\mathbb R}$ est compact,
i.e. $V'(\mathbb R)$ compact, considérant $V'(\mathbb R)$ comme quotient de
$\mathbb R_+^{*(r_1+r_2-1)}\times(\mathbb Z/2\mathbb Z)^{r_1}\times U^{r_2}$,
on voit que ce sera un quotient de de $(\mathbb Z/2\mathbb Z)^{r_1}\times U^{r_2}$,
et l'image de $E_1$ dans ce dernier est donc finie [l'image}

$(\mathrm{iii})\Longrightarrow(\mathrm{iv})$ Si $V'_K$ est un tore compact, le groupe
des car. à \deleted{\unsure{coeff}} valeurs dans \deleted{$\unsure{\widetilde K}$}
$\overline{\mathbb Q}\subset\mathbb C$, regardé comme module sous
$G=\operatorname{Gal}(\widetilde K/\unsure{K|\mathbb Q})$
($\widetilde K\subset\overline{\mathbb Q}$ une ext. de $K$ finie galoisienne sur $\mathbb Q$)
est telle que $\tau\in G$ (la conjugaison complexe) y opère par $-\operatorname{id}_M$,
\inserted{donc est central, donc $M$ est un $G'=G/\unsure{P}$-module, où $G'$}
\deleted{Donc si $G'=G/\unsure{H}$ est le quotient de $G$ qui opère fidèlement sur $M$,}
est le grpe de Galois de la plus gde ss-ext. Weilienne
$K'\subset\widetilde K$ \deleted{\unsure{a défini une extension}}
de $K$ de $\widetilde K$.
\inserted{Prouvons d'abord que $R\supset\operatorname{Ker}(N_{K/L}|V_K)$
i.e. que $V'$ est un \uline{quotient} de $V_L$ considéré comme quotient de $V_K$.}
\deleted{Alors $\unsure{K'}$ \unsure{est} $\tau'=\operatorname{Im}(\tau\mapsto G')$
est dans le centre de $G'$, donc $K'$ est Weilienne.}

\margin{\diagram{\begin{tikzcd}
\deleted{\unsure{\widetilde K}} \arrow[d,no head] & \\
K \arrow[r,no head,"Q"] \arrow[d,no head,"H"'] & \widetilde K \arrow[d,no head,"P"] \arrow[ddl,no head,bend left=35,"G"] \\
L=K\cap K' \arrow[r,no head] \arrow[d,no head] & K' \arrow[dl,no head,"G'=G/P"] \\
\mathbb Q &
\end{tikzcd}}}

$I_K$ \inserted{$=G/Q$} espace homogène sous $G$ des hom de $K$ dans $\overline{\mathbb Q}$

$\mathbb Z[I_K]$ [groupe libre engendré par $I$] groupe des caractères de $T_K$
à valeurs dans $\overline{\mathbb Q}$, comme $G$-module.

\deleted{$\unsure{M}\subset\mathbb Z[\unsure{I}]$ tel que}

$\mathbb Z\subset\mathbb Z[I_K]$ fonctions constantes sur $I$

\diagram{\begin{tikzcd}
\mathbb Z \arrow[r,no head] & N \arrow[r,no head] & \mathbb Z[I_K] \\
\mathbb Z \arrow[u,dotted,"\simeq"] \arrow[rr,no head] & & \mathbb Z(I_L) \arrow[u,hook]
\end{tikzcd}}

$N$ sous $G$-module, tel que $\tau$ opère sur $N/\mathbb Z$ par $-\mathrm{id}$,
donc $\unsure{P|A}$ y opère trivialmt; donc $H$ opère trivialmt sur $N$, donc
$N\subset\mathbb Z[I_K]^P=\mathbb Z[P\backslash I_K]$
\deleted{$\unsure{\mathbb Z}$}.
Or $P\backslash I_K=I_L$, donc on trouve \deleted{$\unsure{\subset}$}
$N\subset\mathbb Z[I_L]$, ce qui signifie que le quotient $V'$ de $V_K$ est majoré par $V_L$.

\margin{$I_K=G/Q$\par
$I_L=G/PQ=P\backslash I_K\deleted{\unsure{Q}}$\par
$\tau\in G$\par d'ordre $2$}
\margin{142}
